% --- File: appendice-riferimento-rapido.tex ---
\chapter{Quick Reference: Metodi di uso frequente (bozza da sistemare)}

In questa appendice sono raccolti i metodi delle librerie standard di Java 21 necessari per risolvere la maggior parte degli esercizi algoritmici senza consultare la documentazione esterna.

\section{Classe String}
Le stringhe sono immutabili. Ogni metodo restituisce una nuova istanza.

\begin{table}[ht!]
\centering
\begin{tabularx}{\textwidth}{@{}lX@{}}
\toprule
\textbf{Metodo} & \textbf{Descrizione} \\ \midrule
\texttt{length()} & Restituisce il numero di caratteri. \\
\texttt{charAt(int i)} & Carattere alla posizione $i$. \\
\texttt{substring(start, end)} & Sottostringa dall'indice \textit{start} a \textit{end-1}. \\
\texttt{contains(CharSequence)} & Verifica se la sequenza e presente. \\
\texttt{split(String regex)} & Divide la stringa in un array di sottostringhe. \\
\texttt{toLowerCase() / toUpperCase()} & Cambio di \textit{case}. \\
\texttt{strip() / trim()} & Rimuove spazi bianchi (strip e piu moderno/Unicode). \\ \bottomrule
\end{tabularx}
\end{table}

\section{Classe Math}
Metodi statici per operazioni matematiche comuni.

\begin{table}[ht!]
\centering
\begin{tabularx}{\textwidth}{@{}lX@{}}
\toprule
\textbf{Metodo} & \textbf{Esempio / Risultato} \\ \midrule
\texttt{Math.min(a, b) / Math.max(a, b)} & Restituisce il minore o il maggiore. \\
\texttt{Math.abs(x)} & Valore assoluto $|x|$. \\
\texttt{Math.sqrt(x)} & Radice quadrata $\sqrt{x}$. \\
\texttt{Math.pow(base, esp)} & Elevamento a potenza (restituisce \texttt{double}). \\
\texttt{Math.round(x)} & Arrotondamento all'intero piu vicino. \\ \bottomrule
\end{tabularx}
\end{table}

\section{BigInteger e BigDecimal}
Metodi per l'aritmetica ad alta precisione.

\begin{table}[ht!]
\centering
\begin{tabularx}{\textwidth}{@{}lX@{}}
\toprule
\textbf{Operazione} & \textbf{Sintassi Java} \\ \midrule
Somma / Sottrazione & \texttt{a.add(b) / a.subtract(b)} \\
Moltiplicazione / Divisione & \texttt{a.multiply(b) / a.divide(b, roundingMode)} \\
Confronto & \texttt{a.compareTo(b)} ($0$ se uguali, $1$ se $a>b$, $-1$ se $a<b$) \\ \bottomrule
\end{tabularx}
\end{table}

\section{Gestione del Tempo (java.time)}
Ricorda: queste classi sono immutabili. Ogni modifica crea un nuovo oggetto.

\begin{table}[ht!]
\centering
\begin{tabularx}{\textwidth}{@{}lX@{}}
\toprule
\textbf{Metodo} & \textbf{Descrizione / Esempio} \\ \midrule
\texttt{LocalDate.now()} & Data corrente del sistema. \\
\texttt{LocalDate.of(y, m, d)} & Crea una data specifica (es. 2026, 4, 2). \\
\texttt{date.plusDays(n)} & Aggiunge $n$ giorni (restituisce nuovo oggetto). \\
\texttt{date.isBefore(other)} & Verifica se la data precede un'altra. \\
\texttt{ChronoUnit.DAYS.between(a, b)} & Calcola i giorni tra due date. \\ \bottomrule
\end{tabularx}
\end{table}

\section{Parsing e Conversione}
Metodi per trasformare stringhe in tipi di dato e viceversa.

\begin{table}[ht!]
\centering
\begin{tabularx}{\textwidth}{@{}lX@{}}
\toprule
\textbf{Operazione} & \textbf{Sintassi Java} \\ \midrule
Da String a int & \texttt{Integer.parseInt("123")} \\
Da String a double & \texttt{Double.parseDouble("12.5")} \\
Da int/double a String & \texttt{String.valueOf(n)} \\
Format di output & \texttt{System.out.printf("\%d - \%.2f", i, d)} \\ \bottomrule
\end{tabularx}
\end{table}